% Cote 119: offprints, reports, courses, CV and letters (typed and printed part): archive cote 119, site page 243 (PDF page 244).
% typescript. Draft compared: none.
% AI-transcribed by Claude Opus 5.5 (claude-opus-5-5, adaptive thinking, effort high), Message Batches API, batch r119_full (msgbatch_01DwZNm5Y7YsKBWyAN4sFs8w), 2026-11-13; not yet checked by eye.
\scanpage{119/243}
\mspage{3}

fortement il sera aux prises avec les problèmes que lui pose la science qui
est en train de se faire, plus grand sera l'effort qu'il devra faire pour
s'arracher à ces problèmes, et la difficulté qu'il éprouvera pour se mettre
dans l'état d'esprit complètement différent que suppose un enseignement à
des étudiants qui ne s'intéressent nullement à cette science comme telle.
D'une part un tel enseignement a de fortes chances d'être faussé par les
exigences théoriques naturellement él\unsure{\over{e}{è}}vées du maître vis-à-vis de son propre
cours, lesquelles ne répondent nullement aux besoins de ses auditeurs,- d'où
gaspillage de temps et d'énergie intellectuelle de ceux-ci. D'autre part, il
y a là un regrettable gaspillage d'énergie intellectuelle du maître également,
dont la fonction enseignante constitue dans le cas envisagé une entrave, et
non un stimulant ou un compl\unsure{\over{e}{é}}ment, à son activité créatrice.

Pour des raisons numériques également, on se trouve devant des
situations paradoxales et même absurdes tant qu'on continue à confondre
les rôles enseignant et créateur des maîtres à l'Université. Il est notoire,
pour toute personne bien informée de la question, que par suite de l'afflux
des étudiants et de la nécessité des enseignements divers à assurer, le
nombre des scientifiques qui sont capables d'un travail scientifique origi-%
nal est de loin inférieur au nombre d'enseignants nécessaires à l'Université.
Ceci, joint au\deleted{x} manque de crédits, a entrainé toute une série de graves
inconvénients qui sont allés en s'aggravant au fil des ann\unsure{\over{e}{é}}es :

a) Les anciens critères pour le recrutement des maîtres par leur
créativité scientifique étant théoriquement maintenus, un nombre relati-%
vement petit et tout à fait insuffisant de maîtres se trouve être effectivement
recruté.

b) Pour pallier \deleted{à} l'insuffisance croissante des effectifs de
maîtres de conférence et de professeurs, ceux-ci sont assistés dans leur
tâche par un nombre également croissant d'assistants et de maîtres assistants,
dont la position par rapport aux maîtres reste subalterne et la promotion
professionnelle problématique, par suite des critères mentionnés dans a).

c) Malgré l'appel aux forces enseignantes auxiliaires (mentionnées
dans b)) l'insuffisance des effectifs enseignants se fait durement sentir aux

\margin{243}
